% =========================================================================
% Glossary (unnumbered, front matter, roman pages). Two pages maximum.
% =========================================================================

\section*{Glossary}
\addcontentsline{toc}{section}{Glossary}

\begingroup\footnotesize\singlespacing   % scoped: an unscoped size switch here shrank the whole body once
Terms specific to this thesis, or given a specific sense here; ordinary
machine-learning vocabulary is not repeated. Chapter and section numbers
point to a term's fullest treatment, not necessarily its first use.

\begin{description}[leftmargin=1.8em,labelsep=0.4em,itemsep=0.5pt,parsep=0pt,topsep=2pt]

\item[ancestor cone (cone)] The upstream subgraph of recorded $U$ that
could have shaped a failure; observation-limited, not a counterfactual or
a causal attribution.

\item[archaeology] The method of re-reading this project's own earlier
campaign archives to check whether a candidate's ``discovery'' already
existed there. A post-hoc descriptive reading, not an endpoint of H1.

\item[audited scores] Scores recomputed from \texttt{task\_history} under
pass@1 on the 100-task bed, last row per (round, task), carried rows
included, rather than the loop's own \texttt{curves.json}, which is not
comparable across arms (\S\ref{sec:evaluator-bias}, \F{46}).

\item[band (same-config band)] The observed range of round-to-round score
movement, $-8\ldots{+}5$ tasks, is the union over the three no-graph seeds'
twenty ship-free same-configuration windows (\S\ref{sec:power}). The range is a descriptive
reference, not a calibrated detection threshold; a move inside it is not, by
itself, evidence of an edit effect.

\item[capability vs efficacy] The distinction, drawn out in the clinic
(\S\ref{sec:clinic}), between whether the harness \emph{can} take an
action (capability) and whether taking it \emph{changes an outcome}
(efficacy). The \emph{capability wall} is the rival account that the
action space, not the readout, cannot reach the relevant lesion.

\item[candidate] A proposed edit to the harness's configuration, tools,
prompt or processors, authored by the Evolver, before it faces the gates.

\item[carried / fresh] A task-round is \emph{fresh} when the task was
re-executed in that round and \emph{carried} when the previous round's
outcome was held fixed by the audit-batch mechanism
(Appendix~\ref{app:cost}); carried rows occur on the graph arm only \F{58}.

\item[the clinic] The ledger's name (F26--F40) for the capability-vs-efficacy
probe regime of \S\ref{sec:clinic}; in that vocabulary a candidate under
test is a \emph{drug} and a same-window paired run a \emph{trial}.

\item[composition graph ($G$)] A typed, static representation of the
exported processor composition: nodes are exported hook points, processors
and slots. Bundles, skills and tools remain schema categories; runtime tool
and nested-agent events appear in $U$.

\item[Critic] The meta-role that judges each candidate, may interrogate
the Evolver, and audits the candidate set against the scoreboard before
the gates run.

\item[Digester] The meta-role that reads every trajectory, pass and fail
alike, and writes a per-task dossier.

\item[dossier] The Digester's per-task write-up of a trajectory: what
happened, and the reusable strategy or fragility it names.

\item[estimand / identified] The estimand is the quantity a design is built
to estimate; a design \emph{identifies} it when its recorded contrast could
pin that quantity down under its assumptions with unlimited data. Two are
distinguished here: the graph-evidence component effect, which the
whole-stack design does not identify, and the whole-stack contrast between
the two packages, which the design targets but the recorded rows do not
deliver as a common-endpoint contrast (Table~\ref{tab:results}).

\item[Evolver] The meta-role that authors zero or more candidate edits,
each as a manifest plus an applied configuration.

\item[\texttt{\textbackslash F\{n\}} / \cls{A}/\cls{B}/\cls{C}] The anchor
notation: every number in this thesis carries an \F{}~tag into the
numbered row of Appendix~\ref{app:ledger} that supports it; the row is
classed \cls{A} (recomputable by a named script), \cls{B} (recorded, not
regenerable) or \cls{C} (judgment).

\item[flip / flip rate] A task whose pass/fail outcome differs between two
rounds run under an identical configuration; the flip rate is the fraction
of the bed that does this (\S\ref{sec:power}).

\item[four layers of mortality] Four coexisting candidate explanations
for why a prescribed edit may fail to produce a lasting measured gain
(Chapter~\ref{ch:results}).

\item[genotype / deployment / phenotype] Three nested identity hashes for
the graph-represented processor surface: its structural configuration; that
configuration plus its runtime overlay; and the two plus the execution edges
actually observed. Tool-registry content is outside these hashes.

\item[harness] Everything around the model that turns it into an agent:
tools, prompts, processors and the control flow that calls them. ``An LLM
agent is a model plus a harness.''

\item[leak-route task] A swing task more than half of whose passes touch
benchmark artifacts, with reference answers appearing verbatim inside
solver trajectories (\S\ref{sec:measurement}, \F{31}).

\item[level 2] Three unrelated senses collide in this thesis: GAIA
difficulty tier~2 (a benchmark property, always written in that form);
rung~R2 of the readout ladder (``fired, and on what''); and the official
cheatsheet's Level-1/Level-2 \emph{capability evidence} that every
candidate must carry (the artifact loads; the mechanism round-trips),
which Appendix~\ref{app:deviations} calls capability evidence throughout.

\item[lift] A localization metric: the hit rate on a shipped candidate's
named target tasks, divided by the base rate of those tasks flipping
anyway. Defined within one arm over that arm's selected shipped
candidates; the two arms' values are not a contrast \F{9a}--\F{9c}.

\item[no-pixel bed] The 100-task GAIA text-only bed used for every number
in this thesis.

\item[pass@1] One scored attempt per task per round; every score in this
thesis is pass@1, and passes from mixed-$k$ rounds do not count
(\S\ref{sec:evaluator-bias}, Appendix~\ref{app:deviations}).

\item[Planner] The meta-role that synthesizes every dossier in a round
into one landscape document for the Evolver.

\item[mechanism-consistent guard event] A post-hoc adjacent-round observation
in which starvation exits and score moved together after concurrent ships,
the legible one a step-countdown processor; no same-window counterfactual
isolates its efficacy, and it was not retained (F39).

\item[processor] A pluggable unit of harness behaviour, composed into
bundles and hung on hook points; the object the composition graph types.

\item[ratchet] The paper's specified --- and, in the released code,
unimplemented --- rule that a task solved in one round may never be given
back in a later one.

\item[readout / readout ladder (R0--R3)] The evidence available about a
candidate edit, in four rungs: \textbf{R0} deployed, \textbf{R1} would have
executed on the failures it cites, \textbf{R2} fired (and on what),
\textbf{R3} changed outcomes. The current implementation assesses R1
against one selected recorded graph; without a replay graph, that check is
recorded as unchecked. These readout rungs are distinct from campaign
rounds R0--R15 and from the GHX switch ladder, numbered L0--L5 in the run
archives. ``Readout'' is this evidence generally.

\item[recorder] The write-only component that logs the composition graph
and $U$ during a run; it decides nothing itself, and what it writes is what
the injection seams and the sixth gate read.

\item[resolution] The model-conditional minimum detectable score change
(\S\ref{sec:resolution}); not a cutoff defined by the observed band, GAIA's
integer scoring resolution, or the VIM instrumentation sense.

\item[same-configuration window] A run, or pair of runs, executed under
one unchanged configuration, used to measure noise rather than an edit's
effect.

\item[seam] A point in the official loop's six stages at which GHX
attaches from outside, without altering what runs there by default.

\item[ship / \texttt{no\_op}] The two terminal decisions a round can
reach: land one or more candidates (\emph{ship}), or change nothing
(\emph{no\_op}).

\item[sixth gate] GHX's added acceptance gate: when a replay graph of one
selected failure the candidate cites exists, its claimed target must appear
(a processor: appear and intervene) in that graph, or the candidate is
refused. Without a replay graph, the gate records an unchecked pass
(\texttt{checked=False}).

\item[survivor (the survivor)] The one positive channel in this study ---
episodic notes prepended to the prompt --- that outlasted every governance
and dormancy failure mode the rest of the loop suffered.

\item[treadmill (re-prescription treadmill)] Independent re-derivation of
the same edit family across clean-template campaigns without cross-campaign
transfer; the pattern does not establish efficacy or forgetting.

\item[typed candidate surface] Typed candidate operations comprising
processor-graph edits and separate tool-registry edits that a build-time
validator can check.

\item[unfolded execution graph ($U$)] The per-invocation DAG of what a run
actually executed, with edges observed rather than reconstructed or
inferred.

\item[whitelist / whitelist campaign] The evidence-admissibility ruling
(2026-08-26) restricting every claim in this thesis to specific campaigns;
a \emph{whitelist campaign} is one admitted under it; the singular phrase \emph{the whitelist graph campaign} denotes graph seed 1.

\end{description}
\endgroup
